\movedhead{Statistical analysis}
\label{sec:statanalysis}\label{app:statanalysis}
The unit of analysis is the matched pair. Trials are averaged within a
condition, pairs give paired differences, 95\% intervals come from a
bootstrap over pairs with 10{,}000 resamples, and two sided $p$ values come
from a sign flipping permutation test with 10{,}000 permutations. Both
dependent variables pass through identical machinery, so the categorical and
continuous estimates are directly comparable in everything but units.

An effect is called non-zero when the 95\% interval excludes zero. That rule,
and not the $p$ value, is the pre-registered decision rule, which matters here
for a reason worth stating plainly. A sign flipping test draws its resolution
from the pairs that moved: flipping the sign of a pair whose difference is
zero changes nothing, so with $n_z$ moving pairs the smallest attainable two
sided $p$ is about $2/2^{n_z}$. At twelve pairs per family a family where five
pairs move has a floor of 0.0625, and no effect size can push it lower. In
most cells the floor sits above 0.05, so the permutation test cannot express
significance there however large the effect, and Holm correction across nine
families cannot be cleared by construction. We therefore report every $p$
alongside its floor and its count of moving pairs, so that evidence which is
weak can be told apart from evidence at the instrument's resolution, and we
rest the non-zero calls on the intervals, which the tie structure does not
affect.

Holm correction is applied across all nine families within each
model and effort cell, reported alongside uncorrected values. It is applied
separately to each measure. Adding a second measure doubles the number of
tests but does not change what a family is, so correcting across both at once
would penalise a family for the existence of a second measure of the same
quantity. Flip rate and within-item entropy describe the categorical answer
and so are reported once per cell.
